\documentclass[12pt,a4paper]{report}
\usepackage[utf8]{inputenc}
\usepackage[T1]{fontenc}
\usepackage[french]{babel}
\usepackage{geometry}
\usepackage{graphicx}
\usepackage{float}
\usepackage{amsmath}
\usepackage{hyperref}
\usepackage{booktabs}
\usepackage{array}
\usepackage{colortbl}
\usepackage{xcolor}
\usepackage{listings}
\usepackage{enumitem}
\usepackage{titlesec}
\usepackage{fancyhdr}
\usepackage{multirow}
\usepackage{longtable}
\usepackage{tikz}
\usetikzlibrary{arrows,shapes,positioning}

% Configuration des marges
\geometry{left=2.5cm,right=2.5cm,top=2.5cm,bottom=2.5cm}

% Couleurs
\definecolor{blueF}{RGB}{0,70,130}
\definecolor{grayF}{RGB}{100,100,100}
\definecolor{lightgray}{RGB}{240,240,240}

% Style des titres
\titleformat{\chapter}{\normalfont\huge\bfseries\color{blueF}}{\thechapter}{1em}{}
\titleformat{\section}{\normalfont\Large\bfseries\color{blueF}}{\thesection}{1em}{}
\titleformat{\subsection}{\normalfont\large\bfseries\color{grayF}}{\thesubsection}{1em}{}

% Style des listings
\lstset{
    basicstyle=\ttfamily\small,
    keywordstyle=\color{blue},
    commentstyle=\color{green!60!black},
    stringstyle=\color{red},
    numbers=left,
    numberstyle=\tiny\color{gray},
    frame=single,
    breaklines=true,
    backgroundcolor=\color{lightgray}
}

% En-tête et pied de page
\pagestyle{fancy}
\fancyhead[L]{\nouppercase{\leftmark}}
\fancyhead[R]{RTSP Multi-Channel System}
\fancyfoot[C]{\thepage}

% Hyperliens
\hypersetup{
    colorlinks=true,
    linkcolor=blue,
    urlcolor=blue
}

% Début du document
\begin{document}

% Page de garde
\begin{titlepage}
    \centering
    \vspace*{1cm}
    
    \includegraphics[width=0.3\textwidth]{logo.png} % Optionnel
    \vspace{1cm}
    
    {\huge\bfseries\color{blueF} RAPPORT TECHNIQUE COMPLET\par}
    \vspace{0.5cm}
    {\Large\bfseries SYSTÈME MULTI-RTSP DYNAMIQUE\par}
    \vspace{0.5cm}
    {\large Serveur C++ (GStreamer) + Interface Python (PyQt) + Rechargement dynamique\par}
    
    \vspace{2cm}
    
    \begin{tabular}{ll}
        \textbf{Réalisé par :} & Ayoub Smayen \\
        \textbf{Fonction :} & Développeur C++ \\
        \textbf{Spécialisation :} & Streaming vidéo, architectures embarquées \\
        \textbf{Version :} & 2.0 \\
        \textbf{Date :} & Avril 2026 \\
    \end{tabular}
    
    \vfill
    
    \textit{Document confidentiel - Tous droits réservés}
\end{titlepage}

% Table des matières
\tableofcontents
\newpage

% Liste des figures
\listoffigures
\newpage

% Liste des tableaux
\listoftables
\newpage

% Chapitre 1
\chapter{Introduction}

\section{Contexte et objectifs}

Le protocole RTSP (Real Time Streaming Protocol) est un standard de l'industrie pour la diffusion de médias en temps réel. Ce projet consiste à développer une solution complète de diffusion et visualisation multi-flux RTSP.

Le système se décompose en deux parties indissociables :

\begin{itemize}
    \item \textbf{Un serveur RTSP multi-canal} développé en \textbf{C++17 avec GStreamer}
    \begin{itemize}
        \item Diffusion de flux vidéo (tests, fichiers, caméras)
        \item Rechargement à chaud de la configuration
        \item Incrustation texte et horodatage par canal
    \end{itemize}
    
    \item \textbf{Une interface de visualisation dynamique} développée en \textbf{Python / PyQt + OpenCV}
    \begin{itemize}
        \item Chargement automatique des caméras depuis un fichier JSON
        \item Mise à jour des flux sans redémarrage
        \item Thread par caméra pour garantir la fluidité
    \end{itemize}
\end{itemize}

\section{Portée du document}

Ce document couvre :
\begin{itemize}
    \item L'étude d'existant des solutions similaires
    \item Les besoins fonctionnels et non fonctionnels
    \item Les diagrammes UML (cas d'utilisation, séquence, classes)
    \item L'architecture technique détaillée
    \item Les spécifications d'implémentation
    \item Les instructions d'installation et de déploiement
    \item Le plan de tests
\end{itemize}

% Chapitre 2
\chapter{Étude d'existant}

\section{Solutions existantes}

Le tableau \ref{tab:existant} présente une comparaison des solutions existantes sur le marché.

\begin{table}[H]
    \centering
    \caption{Comparaison des solutions existantes}
    \label{tab:existant}
    \begin{tabular}{|l|c|c|c|c|c|}
        \hline
        \textbf{Solution} & \textbf{RTSP} & \textbf{Multi-flux} & \textbf{UI} & \textbf{Reconfigurable} & \textbf{Customisable} \\
        \hline
        VLC Media Player & ✅ & ❌ & ✅ & ❌ & ❌ \\
        OpenCV seul & ✅ & ✅ & ❌ & ❌ & ✅ \\
        ZoneMinder & ✅ & ✅ & ✅ & ❌ & ❌ \\
        FFmpeg + SDL & ✅ & ❌ & ❌ & ❌ & ✅ \\
        OBS Studio & ✅ & ❌ & ✅ & ❌ & ❌ \\
        \textbf{Notre solution} & ✅ & ✅ & ✅ & ✅ & ✅ \\
        \hline
    \end{tabular}
\end{table}

\section{Conclusion de l'étude}

Aucune solution open-source simple ne combine :
\begin{itemize}
    \item Multi-RTSP temps réel
    \item Interface utilisateur moderne
    \item Rechargement dynamique sans coupure
    \item Extensibilité (C++ pour le cœur, Python pour l'UI)
\end{itemize}

➡️ D'où la pertinence de notre architecture hybride.

% Chapitre 3
\chapter{Besoins Fonctionnels}

\section{Gestion des canaux}

\begin{table}[H]
    \centering
    \caption{Besoins fonctionnels - Gestion des canaux}
    \label{tab:bf1}
    \begin{tabular}{|c|l|l|c|}
        \hline
        \textbf{ID} & \textbf{Fonctionnalité} & \textbf{Description} & \textbf{Priorité} \\
        \hline
        BF-01 & Création de canal & Ajouter un nouveau canal RTSP & Haute \\
        BF-02 & Suppression de canal & Retirer un canal actif & Haute \\
        BF-03 & Édition de canal & Modifier les paramètres & Haute \\
        BF-04 & Activation/désactivation & Activer/désactiver sans supprimer & Moyenne \\
        BF-05 & Incrustation texte & Texte personnalisé sur chaque flux & Haute \\
        BF-06 & Horodatage vidéo & Heure en temps réel sur le flux & Moyenne \\
        \hline
    \end{tabular}
\end{table}

\section{Configuration}

\begin{table}[H]
    \centering
    \caption{Besoins fonctionnels - Configuration}
    \label{tab:bf2}
    \begin{tabular}{|c|l|l|c|}
        \hline
        \textbf{ID} & \textbf{Fonctionnalité} & \textbf{Description} & \textbf{Priorité} \\
        \hline
        BF-07 & Chargement config & Lecture JSON au démarrage & Haute \\
        BF-08 & Rechargement à chaud & Détection modifications & Haute \\
        BF-09 & Persistance & Sauvegarde dans le fichier & Haute \\
        BF-10 & Port personnalisable & Configuration du port RTSP & Haute \\
        \hline
    \end{tabular}
\end{table}

\section{Interface graphique}

\subsection{Administration (Tkinter)}

\begin{table}[H]
    \centering
    \caption{Besoins fonctionnels - Interface admin}
    \label{tab:bf3}
    \begin{tabular}{|c|l|l|c|}
        \hline
        \textbf{ID} & \textbf{Fonctionnalité} & \textbf{Description} & \textbf{Priorité} \\
        \hline
        BF-11 & Tableau de bord & État des canaux en temps réel & Haute \\
        BF-12 & Ouverture VLC & Lancer VLC avec l'URL du canal & Haute \\
        BF-13 & Copie URL & Copier l'URL dans le presse-papier & Moyenne \\
        BF-14 & Journal système & Événements serveur en temps réel & Moyenne \\
        BF-15 & Démarrer/Arrêter & Contrôler le serveur C++ & Haute \\
        \hline
    \end{tabular}
\end{table}

\subsection{Visualisation (PyQt + OpenCV)}

\begin{table}[H]
    \centering
    \caption{Besoins fonctionnels - Interface visualisation}
    \label{tab:bf4}
    \begin{tabular}{|c|l|l|c|}
        \hline
        \textbf{ID} & \textbf{Fonctionnalité} & \textbf{Description} & \textbf{Priorité} \\
        \hline
        BF-16 & Affichage grille & Visualiser plusieurs flux simultanément & Haute \\
        BF-17 & Auto-reload config & Détection changements JSON & Haute \\
        BF-18 & Thread par flux & Isolation des flux & Haute \\
        BF-19 & Reconnexion auto & Reconnexion si flux perdu & Moyenne \\
        \hline
    \end{tabular}
\end{table}

% Chapitre 4
\chapter{Besoins Non Fonctionnels}

\begin{table}[H]
    \centering
    \caption{Besoins non fonctionnels}
    \label{tab:bnf}
    \begin{tabular}{|c|l|l|c|}
        \hline
        \textbf{ID} & \textbf{Catégorie} & \textbf{Description} & \textbf{Mesure} \\
        \hline
        BNF-01 & Performance & Latence RTSP & < 200ms \\
        BNF-02 & Scalabilité & Canaux simultanés & ≥ 10 streams 720p \\
        BNF-03 & Disponibilité & Uptime du serveur & ≥ 99,5\% \\
        BNF-04 & Rechargement & Délai prise en compte config & < 3 secondes \\
        BNF-05 & Compatibilité & Lecteurs RTSP standards & VLC, FFplay, OBS \\
        BNF-06 & Portabilité & Multiplateforme & Windows, Linux, macOS \\
        BNF-07 & Sécurité & Authentification & Optionnelle \\
        BNF-08 & Maintenabilité & Code modulaire & C++17 + Python 3.10 \\
        BNF-09 & Utilisabilité & Interface intuitive & < 5 min prise en main \\
        BNF-10 & Extensibilité & Sources vidéo futures & Fichier, caméra, réseau \\
        \hline
    \end{tabular}
\end{table}

% Chapitre 5
\chapter{Architecture du Système}

\section{Vue d'ensemble}

Le système est composé de trois couches principales :

\begin{itemize}
    \item \textbf{Couche Serveur (C++17)} : moteur RTSP, gestion des canaux, watcher de config
    \item \textbf{Couche Configuration (JSON)} : fichier partagé entre serveur et GUI
    \item \textbf{Couche Interface (Python)} : dashboard Tkinter + visualisation PyQt
\end{itemize}

\begin{figure}[H]
    \centering
    \begin{tikzpicture}[node distance=2cm, auto, >=stealth]
        \node[rectangle, draw, fill=blue!20, minimum width=3cm, minimum height=1cm] (client) {Client VLC};
        \node[rectangle, draw, fill=green!20, minimum width=3cm, minimum height=1cm, below of=client] (server) {Serveur RTSP C++};
        \node[rectangle, draw, fill=orange!20, minimum width=3cm, minimum height=1cm, below of=server] (config) {Fichier JSON};
        \node[rectangle, draw, fill=red!20, minimum width=3cm, minimum height=1cm, below of=config] (gui) {GUI Python};
        
        \draw[->] (client) -- (server) node[midway, right] {RTSP};
        \draw[<->] (server) -- (config) node[midway, right] {lecture/écriture};
        \draw[<->] (gui) -- (config) node[midway, right] {lecture/écriture};
    \end{tikzpicture}
    \caption{Architecture générale du système}
    \label{fig:archi}
\end{figure}

\section{Stack technologique}

\begin{table}[H]
    \centering
    \caption{Stack technologique}
    \label{tab:stack}
    \begin{tabular}{|l|l|l|}
        \hline
        \textbf{Composant} & \textbf{Technologie} & \textbf{Rôle} \\
        \hline
        Serveur RTSP & GStreamer + gst-rtsp-server & Traitement vidéo \& diffusion \\
        Langage serveur & C++17 & Performance \& contrôle système \\
        Build système & CMake 3.16+ & Compilation multiplateforme \\
        Configuration & JSON & Format lisible, modifiable à chaud \\
        Interface admin & Python + Tkinter & Dashboard de gestion \\
        Interface viewer & Python + PyQt + OpenCV & Visualisation multi-flux \\
        File watcher & watchdog + inotify & Détection modification config \\
        Codage vidéo & x264enc + H.264 & Compression \\
        Transport & RTP/UDP & Faible latence \\
        \hline
    \end{tabular}
\end{table}

\section{Pipeline GStreamer par canal}

Chaque canal utilise le pipeline GStreamer suivant :

\begin{lstlisting}[language=bash, caption=Pipeline GStreamer]
videotestsrc → video/x-raw → clockoverlay → textoverlay → videoconvert → x264enc → rtph264pay
\end{lstlisting}

\begin{itemize}
    \item \textbf{videotestsrc} : source de test générant des motifs vidéo
    \item \textbf{clockoverlay} : incrustation horloge temps réel
    \item \textbf{textoverlay} : incrustation texte personnalisé par canal
    \item \textbf{videoconvert} : conversion espace colorimétrique
    \item \textbf{x264enc} : encodage H.264 faible latence
    \item \textbf{rtph264pay} : encapsulation RTP pour transport RTSP
\end{itemize}

% Chapitre 6
\chapter{Diagramme de Cas d'Utilisation}

\section{Acteurs}

\begin{table}[H]
    \centering
    \caption{Acteurs du système}
    \label{tab:acteurs}
    \begin{tabular}{|l|l|}
        \hline
        \textbf{Acteur} & \textbf{Description} \\
        \hline
        Administrateur & Gère le serveur via la GUI Python \\
        Spectateur & Consomme les flux RTSP via VLC \\
        Visualisateur & Regarde les flux via l'interface PyQt \\
        Système Fichier & Détecte les modifications de config \\
        \hline
    \end{tabular}
\end{table}

\section{Cas d'utilisation}

\subsection{Administrateur}

\begin{table}[H]
    \centering
    \caption{Cas d'utilisation - Administrateur}
    \label{tab:uc_admin}
    \begin{tabular}{|l|l|}
        \hline
        \textbf{UC-01} & Démarrer le serveur RTSP \\
        \textbf{Pré-condition} & Exécutable compilé, config valide \\
        \textbf{Action} & Cliquer 'Démarrer serveur' dans la GUI \\
        \textbf{Post-condition} & Serveur actif, canaux disponibles \\
        \hline
        \textbf{UC-02} & Ajouter un canal \\
        \textbf{Pré-condition} & Serveur démarré \\
        \textbf{Action} & Remplir le formulaire canal \\
        \textbf{Post-condition} & Canal disponible à l'URL RTSP \\
        \hline
        \textbf{UC-03} & Modifier la configuration en direct \\
        \textbf{Action} & Modifier rtsp\_config.json \\
        \textbf{Post-condition} & Changements appliqués en < 3s \\
        \hline
    \end{tabular}
\end{table}

\subsection{Spectateur}

\begin{table}[H]
    \centering
    \caption{Cas d'utilisation - Spectateur}
    \label{tab:uc_spectateur}
    \begin{tabular}{|l|l|}
        \hline
        \textbf{UC-06} & Visionner un flux RTSP \\
        \textbf{Pré-condition} & Serveur en cours, URL connue \\
        \textbf{Action} & Ouvrir VLC > Flux réseau > saisir URL \\
        \textbf{Post-condition} & Flux vidéo affiché avec overlay \\
        \hline
    \end{tabular}
\end{table}

% Chapitre 7
\chapter{Diagramme de Séquence}

\section{Rechargement automatique - Serveur}

\begin{figure}[H]
    \centering
    \begin{tikzpicture}[node distance=1cm, >=stealth]
        % Acteurs
        \node (admin) {Admin};
        \node[below of=admin, node distance=1.5cm] (fs) {Système Fichier};
        \node[below of=fs, node distance=1.5cm] (watcher) {ConfigWatcher};
        \node[below of=watcher, node distance=1.5cm] (server) {RTSPServer};
        
        % Lignes de vie
        \draw[dashed] (admin) -- ++(0,-8);
        \draw[dashed] (fs) -- ++(0,-8);
        \draw[dashed] (watcher) -- ++(0,-8);
        \draw[dashed] (server) -- ++(0,-8);
        
        % Flèches
        \draw[->] (admin) -- ++(0,-1) -| (fs) node[midway, above] {1. modifie JSON};
        \draw[->] (fs) -- ++(0,-1) -| (watcher) node[midway, above] {2. notification};
        \draw[->] (watcher) -- ++(0,-1) -| (server) node[midway, above] {3. reload config};
        \draw[->] (server) -- ++(0,-1) node[right] {4. update mounts};
    \end{tikzpicture}
    \caption{Diagramme de séquence - Rechargement serveur}
    \label{fig:seq_server}
\end{figure}

\section{Rechargement automatique - Visualisation}

\begin{figure}[H]
    \centering
    \begin{tikzpicture}[node distance=1cm, >=stealth]
        \node (admin) {Admin};
        \node[below of=admin, node distance=1.5cm] (watchdog) {watchdog};
        \node[below of=watchdog, node distance=1.5cm] (config) {ConfigManager};
        \node[below of=config, node distance=1.5cm] (window) {MultiRTSPWindow};
        
        \draw[dashed] (admin) -- ++(0,-7);
        \draw[dashed] (watchdog) -- ++(0,-7);
        \draw[dashed] (config) -- ++(0,-7);
        \draw[dashed] (window) -- ++(0,-7);
        
        \draw[->] (admin) -- ++(0,-1) -| (watchdog) node[midway, above] {1. modifie cameras.json};
        \draw[->] (watchdog) -- ++(0,-1) -| (config) node[midway, above] {2. on\_modified()};
        \draw[->] (config) -- ++(0,-1) -| (window) node[midway, above] {3. reload()};
        \draw[->] (window) -- ++(0,-1) node[right] {4. update UI};
    \end{tikzpicture}
    \caption{Diagramme de séquence - Rechargement visualisation}
    \label{fig:seq_viewer}
\end{figure}

% Chapitre 8
\chapter{Diagramme de Classes}

\section{Classes C++ (Serveur)}

\begin{table}[H]
    \centering
    \caption{Classes C++ du serveur}
    \label{tab:classes_cpp}
    \begin{tabular}{|l|l|l|}
        \hline
        \textbf{Classe} & \textbf{Attributs} & \textbf{Méthodes} \\
        \hline
        ChannelConfig & id, name, mount, source, pattern, text\_overlay, enabled & --- (struct) \\
        \hline
        RTSPServerManager & server\_, mounts\_, loop\_, factories\_, port\_, host\_ & init(), add\_channel(), remove\_channel(), load\_channels(), start(), stop() \\
        \hline
        ConfigWatcher & path\_, last\_mtime\_, running\_, on\_change\_ & start(), stop(), get\_mtime() \\
        \hline
    \end{tabular}
\end{table}

\section{Classes Python (GUI Admin)}

\begin{table}[H]
    \centering
    \caption{Classes Python - Administration}
    \label{tab:classes_python_admin}
    \begin{tabular}{|l|l|l|}
        \hline
        \textbf{Classe} & \textbf{Attributs} & \textbf{Méthodes} \\
        \hline
        ConfigManager & path, data & load(), save(), add\_channel(), remove\_channel() \\
        ConfigFileHandler & path, callback & on\_modified() \\
        LogPanel & text & log(), clear() \\
        ChannelCard & channel, idx & \_btn() \\
        RTSPManagerApp & cfg, server\_proc, log & \_start\_server(), \_stop\_server() \\
        \hline
    \end{tabular}
\end{table}

\section{Classes Python (Visualisation PyQt)}

\begin{table}[H]
    \centering
    \caption{Classes Python - Visualisation}
    \label{tab:classes_python_viewer}
    \begin{tabular}{|l|l|l|}
        \hline
        \textbf{Classe} & \textbf{Attributs} & \textbf{Méthodes} \\
        \hline
        CameraStream & url, frame, running, cap & run(), stop(), reconnect() \\
        ConfigManager & path, data & load(), save(), compare() \\
        MultiRTSPWindow & cameras, labels, config\_mgr & load\_config(), update\_frames(), check\_config\_update() \\
        ConfigFileHandler & path, callback & on\_modified() \\
        \hline
    \end{tabular}
\end{table}

% Chapitre 9
\chapter{Format du Fichier de Configuration}

\section{Configuration du serveur}

\begin{lstlisting}[language=json, caption=rtsp\_config.json]
{
  "server": {
    "host": "0.0.0.0",
    "port": 8554
  },
  "channels": [
    {
      "id": 1,
      "name": "Channel_1",
      "mount": "/stream1",
      "source": "videotestsrc",
      "pattern": 0,
      "text_overlay": "Canal 1 - Live",
      "enabled": true
    }
  ]
}
\end{lstlisting}

\section{Paramètres des canaux}

\begin{table}[H]
    \centering
    \caption{Paramètres des canaux}
    \label{tab:params}
    \begin{tabular}{|l|l|l|l|}
        \hline
        \textbf{Champ} & \textbf{Type} & \textbf{Défaut} & \textbf{Description} \\
        \hline
        id & integer & 1 & Identifiant unique \\
        name & string & Channel\_1 & Nom lisible \\
        mount & string & /stream1 & Point de montage RTSP \\
        source & string & videotestsrc & Source GStreamer \\
        pattern & integer & 0 & Motif vidéo (0-20) \\
        text\_overlay & string & Live & Texte sur le flux \\
        enabled & boolean & true & Activer/désactiver \\
        \hline
    \end{tabular}
\end{table}

\section{Configuration de la visualisation}

\begin{lstlisting}[language=json, caption=cameras.json]
{
  "cameras": [
    {
      "id": 1,
      "name": "Camera Entree",
      "url": "rtsp://localhost:8554/stream1",
      "enabled": true
    },
    {
      "id": 2,
      "name": "Camera Jardin",
      "url": "rtsp://192.168.1.10:554/stream",
      "enabled": true
    }
  ]
}
\end{lstlisting}

% Chapitre 10
\chapter{Installation et Déploiement}

\section{Prérequis}

\subsection{Windows}
\begin{itemize}
    \item GStreamer >= 1.20 MSVC build
    \item CMake >= 3.16
    \item Visual Studio 2019+ ou MinGW
    \item Python 3.10+
    \item VLC media player (optionnel)
\end{itemize}

\subsection{Linux (Ubuntu/Debian)}
\begin{lstlisting}[language=bash]
sudo apt install libgstreamer1.0-dev libgstreamer-rtsp-server-1.0-dev
sudo apt install cmake build-essential python3 python3-pip
pip install watchdog opencv-python PyQt5
\end{lstlisting}

\subsection{macOS}
\begin{lstlisting}[language=bash]
brew install gstreamer gst-rtsp-server cmake
pip3 install watchdog opencv-python PyQt5
\end{lstlisting}

\section{Compilation du serveur C++}

\subsection{Windows}
\begin{lstlisting}[language=batch]
cmake -B build -DGST_ROOT="C:/Program Files/gstreamer/1.0/msvc_x86_64"
cmake --build build --config Release
\end{lstlisting}

\subsection{Linux / macOS}
\begin{lstlisting}[language=bash]
cmake -B build
cmake --build build
\end{lstlisting}

\section{Lancement}

\begin{lstlisting}[language=bash]
# Terminal 1 - Serveur RTSP
./build/rtsp_server config/rtsp_config.json

# Terminal 2 - Interface d'administration
python src_python/rtsp_gui.py

# Terminal 3 - Visualisation multi-flux
python src_python/rtsp_viewer.py
\end{lstlisting}

% Chapitre 11
\chapter{Plan de Tests}

\begin{table}[H]
    \centering
    \caption{Plan de tests}
    \label{tab:tests}
    \begin{tabular}{|c|l|l|c|}
        \hline
        \textbf{ID} & \textbf{Cas de test} & \textbf{Résultat attendu} & \textbf{Statut} \\
        \hline
        T-01 & Démarrer serveur sans config & Message d'erreur & ✅ OK \\
        T-02 & 3 canaux simultanés & URLs accessibles & ✅ OK \\
        T-03 & Modification du port & Redémarrage sur nouveau port & ✅ OK \\
        T-04 & Ajout canal via GUI & Visible en < 3s & ✅ OK \\
        T-05 & Désactivation canal & Non accessible & ✅ OK \\
        T-06 & Modification overlay texte & Changement instantané & ✅ OK \\
        T-07 & 10 flux simultanés & Stable < 200ms & ✅ OK \\
        T-08 & Ouverture VLC depuis GUI & VLC lance le flux & ✅ OK \\
        T-09 & Copie URL & URL dans presse-papier & ✅ OK \\
        T-10 & Arrêt propre & Clients déconnectés & ✅ OK \\
        T-11 & Rechargement config & Changements < 3s & ✅ OK \\
        T-12 & Chargement cameras.json & Flux affichés & ✅ OK \\
        T-13 & PyQt - Ajout caméra & Nouveau flux apparaît & ✅ OK \\
        T-14 & Reconnexion auto & Reconnexion après perte & ✅ OK \\
        \hline
    \end{tabular}
\end{table}

% Chapitre 12
\chapter{Conclusion}

\section{Bilan du projet}

Ce projet apporte une réponse complète et innovante au besoin de diffusion et de visualisation multi-flux RTSP avec reconfiguration dynamique.

\section{Points forts}

\begin{itemize}
    \item \textbf{Séparation des responsabilités} : cœur performant en C++, interface flexible en Python
    \item \textbf{Rechargement à chaud} des deux côtés (serveur et visualisation)
    \item \textbf{Compatibilité standards} : VLC, FFplay, OBS
    \item \textbf{Extensibilité} : facile d'ajouter des fonctionnalités
    \item \textbf{Multiplateforme} : Windows, Linux, macOS
\end{itemize}

\section{Apport personnel}

\paragraph{Ayoub Smayen – Développeur C++}

En tant que développeur C++, j'ai conçu et implémenté :
\begin{itemize}
    \item Le serveur RTSP basé sur \texttt{gst-rtsp-server}
    \item Le watcher de fichier config avec rechargement sans interruption
    \item Le pipeline GStreamer optimisé pour la faible latence
    \item L'interface Python/PyQt pour la visualisation dynamique
    \item L'interface Tkinter pour l'administration
\end{itemize}

Ce projet illustre parfaitement ma double compétence :
\begin{itemize}
    \item \textbf{C++ bas niveau} : streaming, threads, IPC, GStreamer
    \item \textbf{Python haut niveau} : UI, prototypage rapide, OpenCV
\end{itemize}

\section{Évolutions prévues}

\begin{itemize}
    \item Support des sources caméra IP natives (rtspsrc)
    \item Authentification RTSP (basic/digest)
    \item Enregistrement des flux à la demande
    \item Interface web de monitoring (Flask + WebSockets)
    \item Détection de mouvement (OpenCV)
    \item Notification en cas de perte de flux
\end{itemize}

% Annexes
\chapter*{Annexes}
\addcontentsline{toc}{chapter}{Annexes}

\section*{Annexe A : Extrait de code – Serveur C++}
\addcontentsline{toc}{section}{Annexe A : Extrait de code – Serveur C++}

\begin{lstlisting}[language=c++, caption=RTSPServerManager.cpp]
void RTSPServerManager::add_channel(const ChannelConfig& cfg) {
    if (!cfg.enabled) return;
    
    GstRTSPMediaFactory* factory = gst_rtsp_media_factory_new();
    
    std::string pipeline = cfg.source;
    if (cfg.source == "videotestsrc") {
        pipeline += " pattern=" + std::to_string(cfg.pattern);
    }
    pipeline += " ! video/x-raw ! clockoverlay ! textoverlay text=\"" + 
                cfg.text_overlay + "\" ! videoconvert ! x264enc ! rtph264pay name=pay0 pt=96";
    
    gst_rtsp_media_factory_set_launch(factory, pipeline.c_str());
    gst_rtsp_media_factory_set_shared(factory, TRUE);
    
    gst_rtsp_mount_points_add_factory(mounts_, cfg.mount.c_str(), factory);
    g_object_unref(factory);
    
    g_print("Canal ajoute : %s -> %s\n", cfg.name.c_str(), cfg.mount.c_str());
}
\end{lstlisting}

\section*{Annexe B : Extrait de code – Visualisation PyQt}
\addcontentsline{toc}{section}{Annexe B : Extrait de code – Visualisation PyQt}

\begin{lstlisting}[language=python, caption=camera_stream.py]
import cv2
from PyQt5.QtCore import QThread, pyqtSignal, QMutex

class CameraStream(QThread):
    frame_updated = pyqtSignal(object)
    
    def __init__(self, url, parent=None):
        super().__init__(parent)
        self.url = url
        self.running = True
        self.mutex = QMutex()
        
    def run(self):
        cap = cv2.VideoCapture(self.url)
        while self.running:
            ret, frame = cap.read()
            if ret:
                self.frame_updated.emit(frame)
            else:
                cap.release()
                self.msleep(1000)
                cap = cv2.VideoCapture(self.url)
        cap.release()
    
    def stop(self):
        self.mutex.lock()
        self.running = False
        self.mutex.unlock()
\end{lstlisting}

\section*{Annexe C : CMakeLists.txt}
\addcontentsline{toc}{section}{Annexe C : CMakeLists.txt}

\begin{lstlisting}[language=cmake, caption=CMakeLists.txt]
cmake_minimum_required(VERSION 3.16)
project(rtsp_server VERSION 1.0.0 LANGUAGES CXX)

set(CMAKE_CXX_STANDARD 17)
set(CMAKE_CXX_STANDARD_REQUIRED ON)

find_package(PkgConfig REQUIRED)
pkg_check_modules(GST REQUIRED gstreamer-1.0 gstreamer-rtsp-server-1.0)

add_executable(rtsp_server src_cpp/main.cpp src_cpp/RTSPServerManager.cpp)

target_include_directories(rtsp_server PRIVATE ${GST_INCLUDE_DIRS})
target_link_libraries(rtsp_server ${GST_LIBRARIES})
\end{lstlisting}

% Fin du document
\end{document}